\chapter{Endings and Retrospective Meaning}\label{ch:endings}

Alternate endings are the most familiar kind of variation. Films have been released with different endings for different markets, recut after test screenings, and restored in director's versions with their original conclusions. \emph{Clue} (1985) went further, distributing three endings so that different theaters showed different ones, and which ending a viewer saw depended on where they went \cite{clue}. That was the intertheater variation of \Cref{ch:local}, applied to endings. Selective cinema changes one thing about them: the endings are seen at the same time, in the same room, by people who then leave together. This chapter examines what an ending does to the film before it, and what a family of endings does to an audience.

\section{The branch}

The simplest family of endings shares everything until a final branch and then resolves differently. Before the branch, the realizations are identical, or differ only in ways that establish nothing relevant. After it, they establish contradictory things: the character survives in one and dies in the other, the plan succeeds or fails, the suspect is guilty or innocent.

In the terms of \Cref{ch:switch-admissibility}, the branch is the family's last structural rendezvous. Before it, the realization graph is complete. After it, the graph is empty: a viewer who has seen one ending begin cannot move into another without being told something that contradicts what they have just seen. The branch point itself is a moment of choice. A viewer who turns the dial at the branch, and only there, selects an ending without interrupting the film and without any incoherence. Selective cinema therefore provides something branching video provides only by stopping the story: a choice of ending made in the flow of the screening, by each viewer independently, with nobody else in the room affected.

In image distance the branch family is the cheapest of all. Its realizations coincide until the branch, so adaptive gating keeps the glasses open for most of the film, and the costs of multiplexing are paid only in the last minutes. It is the natural first project for anyone building such a system.

\section{What endings can be}

An ending cannot be anything at all. It must be consistent with what the film has already established.

\begin{proposition}
Let a family of endings share everything before a branch at time $t_b$, and let $E_{\mathrm{sh}}=\mathrm{Est}(F^{\le t_b})\cap\mathcal R$ be what the shared portion establishes that any ending depends on. Every coherent ending commits only to literals consistent with $E_{\mathrm{sh}}$.
\end{proposition}

\begin{proof}
Each realization must admit its own continuation. An ending committing to the negation of a literal in $E_{\mathrm{sh}}$ would fail consistency against its own history.
\end{proof}

The proposition is simple and its consequence is not. The space of possible endings is bounded by what the shared portion leaves \emph{unestablished}. Everything the film settles before the branch is common to every ending; everything it leaves open is where the endings can differ. A writer's control over a family of endings is therefore, to a large extent, control over what the film declines to settle. A shared portion that establishes too much leaves the endings no room. One that establishes too little leaves the story shapeless. The craft is in presenting a great deal while settling only what every ending can share.

\section{Terminal, seeded, and retrospective}

There are three ways to use that space, and they distribute fact distance differently.

A \emph{terminal} family places all of its divergence after the branch. The shared portion is neutral with respect to the endings; nothing before the branch points toward one rather than another. The fact distance is concentrated in the final minutes, which keeps switching open until then. The risk is that the endings feel arbitrary. An ending the film has not prepared for can feel attached rather than arrived at.

A \emph{seeded} family prepares each ending before the branch. Seeds can be placed in two ways. If each realization carries only the seeds of its own ending, the realizations diverge before the branch, establishing different foreshadowing, and the last structural rendezvous moves earlier. Fact distance is distributed backward through the film, and switching closes sooner. If instead the shared portion carries the seeds of every ending, the realizations stay identical until the branch, but the shared portion must support several futures at once. It becomes deliberately underdetermined, rich in possibilities none of which it confirms. Seeding in this second way is closer to writing a good mystery than to writing a single story, and it is the harder and more interesting form.

A \emph{retrospective} family shares its footage almost entirely and lets the ending change what the footage meant. A final disclosure recasts earlier scenes: a helpful character turns out to have been manipulating everyone, or a threat turns out to have been a misunderstanding. The same shots are seen again in memory, differently. This works only because of the proposition above. The earlier scenes presented the character's helpfulness without establishing it as fact; they established that the character acted helpfully, and left the motive unsettled. Every ending is consistent with what was shown because what was shown was carefully less than what was implied. Retrospective endings live in the gap between presentation and establishment that \Cref{ch:depth} identified from the other side.

\section{An ending is not a scene}

It follows that two realizations with different endings are not identical until the branch, even when their frames are. The same shot of a character's face means one thing in a film that ends in betrayal and another in a film that ends in loyalty. Viewers will not know which they are watching until the end, but afterward they will remember the earlier scenes through the ending they saw. Fact distance, as defined so far, is a property of what the film establishes as it goes. Meaning also has a retrospective component that the definition does not capture: the ending reaches back and changes the significance of everything before it, without changing a single literal the earlier scenes established.

This is why a family of endings is more than a film with a choice at the end. Viewers of different endings have, in a real sense, seen different films throughout, even though they looked at the same frames for most of the running time. The divergence was postponed, not avoided.

\section{Leaving together}

The social consequence is unusual. Two people who watched different endings leave the theater having shared nearly everything and resolved it differently. Their conversation does not begin with what they thought of the ending, as it would after an ordinary film. It begins with which ending each of them saw.

For viewers who did not know the endings differed, that discovery can be disorienting, even unwelcome: the realization that a companion was shown something else. For viewers who knew, it can be the best part of the work. They compare, reconstruct each other's endings, and reinterpret the shared portion in light of both. The film is completed in the conversation. \Cref{ch:after-screening} takes this up and argues that the strongest uses of the form may treat the post-screening conversation as its final act.

The ending is also where the form is most fragile. A leak of a fraction of a second from another stream, which \Cref{ch:phase} showed is hard to eliminate entirely, costs almost nothing in the middle of a genre pair. In the last minute of a family of endings, a glimpse of a face that should be dead can ruin the conclusion. The acceptable level of leakage is not uniform through the film. It should tighten as the ending approaches, and the blanking intervals, light budget, and exhibition settings may need to change with it.

\section{The ending nobody chose}

\Cref{ch:unfiltered} showed that the room without glasses sees the average of the streams. In a family of endings, that composite is quiet for most of the film, because the streams coincide and the unaided viewer sees the film correctly. At the branch it becomes something else: every ending superimposed, the survivor and the corpse at once, the plan succeeding and failing in the same frame.

Ordinarily this would be a defect to minimize. It could also be designed. The composite of all endings is the undecided ending, the room's own version of the conclusion, in which nothing has yet been settled. A viewer who lifts the glasses at the branch sees the story before it chooses. Composing endings so that their average is itself meaningful is difficult, subject to the one-stream cost of the proposition in \Cref{ch:unfiltered}, and possible only where the endings share their staging. But where it can be done, it gives the form an image no other medium has: a single frame containing every way the story could have ended, visible to anyone who looks without choosing.
